\documentclass[12pt]{ximera}
\title{Homework 4: Basic Principles of Probability}
\begin{document}
\begin{abstract}
Sections 7.3--7.4: inclusion-exclusion applied to probabilities, sample spaces and
whether their outcomes are equally likely, and the Monty Hall Problem.
\end{abstract}
\maketitle

\section*{Sections 7.3--7.4: Basic Principles of Probability}

\begin{problem}
Suppose the weather report predicts that there is a $40\%$ chance of rain and a
$35\%$ chance of high winds, with a $15\%$ chance of both. Use inclusion-exclusion
or a Venn diagram to answer the following.

\begin{prompt}
\textbf{(a)} What is the probability that there will be rain or high winds (or
both)?

$\answer[tolerance=0.001]{0.6}$

\begin{hint}
$P(R\cup W)=P(R)+P(W)-P(R\cap W)$. Enter your answer as a decimal.
\end{hint}

\begin{feedback}
\[
P(R\cup W)=P(R)+P(W)-P(R\cap W)=0.40+0.35-0.15=0.60,
\]
so there is a $60\%$ chance of rain or high winds (or both).
\end{feedback}
\end{prompt}

\begin{prompt}
\textbf{(b)} What is the probability that there will be no rain and no high winds?

$\answer[tolerance=0.001]{0.4}$

\begin{hint}
``No rain and no wind'' is the complement of ``rain or wind''.
\end{hint}

\begin{feedback}
``Neither'' is the complement of ``either'', so
\[
P\left(\olsi{R}\cap\olsi{W}\right)=1-P(R\cup W)=1-0.60=0.40,
\]
a $40\%$ chance.
\end{feedback}
\end{prompt}
\end{problem}

\begin{problem}
Suppose we run an experiment where each trial consists of flipping a coin three
times. Let $H$ denote heads and $T$ denote tails.

\begin{prompt}
\textbf{(a)} What is the sample space for this experiment if all we want to record
is how many heads are flipped in each trial?

The sample space has $\answer{4}$ outcomes. Listed in increasing order, they are
$\set{\answer{0},\ \answer{1},\ \answer{2},\ \answer{3}}$.

\begin{feedback}
If we record only the \emph{number} of heads, the possible results are
$0$, $1$, $2$, or $3$ heads, so the sample space is $\set{0,1,2,3}$ --- four
outcomes.
\end{feedback}
\end{prompt}

\begin{prompt}
\textbf{(b)} What is the sample space for the experiment if we also want to record
the order of heads and tails flipped?

This sample space has $\answer{8}$ outcomes.

\begin{hint}
Each of the three flips independently has $2$ possible results.
\end{hint}

\begin{feedback}
Recording order, the sample space is
\[
\set{HHH,\ HHT,\ HTH,\ THH,\ HTT,\ THT,\ TTH,\ TTT},
\]
which has $2^3=8$ outcomes.
\end{feedback}
\end{prompt}

\begin{prompt}
\textbf{(c)} Are the probabilities of the outcomes in part (a) all equal?

\begin{multipleChoice}
\choice{Yes --- there are four outcomes, so each has probability $\frac14$}
\choice[correct]{No --- some counts of heads arise from more ordered outcomes than others}
\end{multipleChoice}

\begin{feedback}
No. The eight \emph{ordered} outcomes in part (b) are equally likely, each with
probability $\frac18$, but they are not spread evenly over the four counts:
\[
P(0)=\tfrac18,\quad P(1)=\tfrac38,\quad P(2)=\tfrac38,\quad P(3)=\tfrac18,
\]
because $1$ head happens in three ways ($HTT$, $THT$, $TTH$) while $0$ heads happens
in only one way ($TTT$).

This is the key lesson: ``there are $4$ outcomes'' does not by itself mean each has
probability $\frac14$. The outcomes must be equally likely first.
\end{feedback}
\end{prompt}
\end{problem}

\begin{problem}
``Probabilities are expressions of our ignorance about the world, and new information
can change the extent of our ignorance.'' --- Sasha Volokh, on the Monty Hall Problem.

There is a famous problem in probability theory called the Monty Hall Problem, named
after the original host of the show \emph{Let's Make A Deal}. It goes as follows.

Suppose you're on a game show, and you're given the choice of three doors. Behind one
door is a car; behind the other two are goats. You pick a door, and then the host, who
knows what's behind all the doors, opens a different door, revealing a goat. He then
says to you, ``Do you want to switch your door choice?'' Is it to your advantage to
switch?

Notice the following facts:
\begin{enumerate}
\item If the contestant's original (random) door choice is correct, then the winning
      decision is to stay.
\item If the contestant's original (random) door choice is not correct, then
      whichever door the host doesn't open is the one with the car, so the winning
      decision is to switch.
\end{enumerate}

\begin{prompt}
\textbf{(a)} What is the probability that the contestant's original (random) choice of
door is correct? Enter a fraction.

$\answer{1/3}$

\begin{feedback}
The car is equally likely to be behind any of the three doors, and the contestant's
first pick is made before any door is opened. So the pick is correct with probability
$\frac13$.
\end{feedback}
\end{prompt}

\begin{prompt}
\textbf{(b)} What is the probability that the contestant's original (random) choice of
door is \emph{not} correct? Enter a fraction.

$\answer{2/3}$

\begin{feedback}
This is the complement of part (a): $1-\frac13=\frac23$.
\end{feedback}
\end{prompt}

\begin{prompt}
\textbf{(c)} Based on your answers to parts (a) and (b), should you switch your door
choice?

\begin{multipleChoice}
\choice[correct]{Yes --- switching wins with probability $\frac23$}
\choice{No --- staying wins with probability $\frac23$}
\choice{It makes no difference --- each remaining door has probability $\frac12$}
\end{multipleChoice}

\begin{feedback}
Switch. By fact 2, switching wins exactly when the original pick was wrong, which
happens with probability $\frac23$. Staying wins only when the original pick was right,
probability $\frac13$.

The tempting answer --- ``two doors are left, so it's $50$--$50$'' --- ignores that the
host's choice is not random: he \emph{always} reveals a goat, and he can never open
your door. Opening a door gives you no new information about your own door, which
stays at $\frac13$, so the remaining $\frac23$ all transfers to the other closed door.
That is the sense of the quotation: the new information changed what you know about
one door, but not the other.
\end{feedback}
\end{prompt}
\end{problem}

\end{document}
